\section{Supporting diagnostics}
The following figures complete the numerical checks described in the main section.
All inputs and results remain synthetic. Every plot has an included EPS source and
a vector PDF companion for direct compilation with a PDF engine.

\begin{figure}[htbp]\centering
\ReportGraphEPS[width=0.88\linewidth]{training_losses.eps}
\caption{Weighted total loss and unweighted data, physics and initial-condition
losses for the selected annual run. The horizontal coordinate combines joint Adam
steps and subsequent L-BFGS closure evaluations. The warmup is omitted, and line-search
trials need not be monotonic.}
\label{fig:yash-losses}
\end{figure}

\begin{figure}[htbp]\centering
\ReportGraphEPS[width=\linewidth]{network_ode_consistency.eps}
\caption{Neural predictions and the ODE trajectory with inferred parameters on
the training interval. Their agreement checks whether fitting the network also
produces a consistent mechanistic trajectory.}
\label{fig:yash-consistency}
\end{figure}

\begin{figure}[htbp]\centering
\ReportGraphEPS[width=0.82\linewidth]{identifiability.eps}
\caption{Singular values of the normalized training-trajectory sensitivity matrix
with respect to log parameters, evaluated at the selected annual estimate. Locally
varying all six parameters introduces less sensitive directions than varying only
$r,a,m$. This is not an all-six PINN fit or a proof of global identifiability.}
\label{fig:yash-identifiability}
\end{figure}

\begin{figure}[htbp]\centering
\ReportGraphEPS[width=\linewidth]{residual_check.eps}
\caption{Hare and lynx residuals evaluated at 307 offset points in the training
interval after fitting. These points diagnose error between collocation points and
do not tune the model. Residual units are inverse years.}
\label{fig:yash-residuals}
\end{figure}
